\section{The lunar datum and its Fisher information}
\label{sec:setup}

\subsection{The seven-parameter datum}

A geodetic datum is the set of frame degrees of freedom that internal geometry
cannot determine. For a rigid body observed by measurements that depend only on
\emph{relative} positions, that set is the seven-parameter similarity (Helmert)
group: three translations $\ttrans=(t_x,t_y,t_z)$, one scale $s$, and three
small rotations $\rot=(\theta_x,\theta_y,\theta_z)$. We collect them in the datum
vector
\begin{equation}
  \datum = (t_x,\,t_y,\,t_z,\,s,\,\theta_x,\,\theta_y,\,\theta_z)^\top \in \R^7,
  \label{eq:datumvec}
\end{equation}
ordered so that indices $0$ and $3$ are the origin-$X$ translation and the scale.
Acting on a body-fixed point $\bm{p}\in\R^3$, the linearized similarity transform is
\begin{equation}
  \bm{p}' \;=\; \bm{p} + \ttrans + s\,\bm{p} + \rot\times\bm{p} + O(\|\datum\|^2),
  \label{eq:helmert}
\end{equation}
whose Jacobian with respect to the datum is the $3\times7$ point-Jacobian
\begin{equation}
  \bm{J}(\bm{p}) \;=\;
  \frac{\partial \bm{p}'}{\partial \datum}
  \;=\;
  \bigl[\; \bm{I}_3 \;\big|\; \bm{p} \;\big|\; -[\bm{p}]_\times \;\bigr],
  \qquad
  [\bm{p}]_\times \rot = \bm{p}\times\rot,
  \label{eq:pointjac}
\end{equation}
in which the first three columns are the translation, the fourth is the scale
(the point's own coordinates), and the last three are the generators of rotation.
The datum is expressed in the lunar principal-axis (PA) body frame in which the
ILRF is defined \citep{sosnica2025ilrf,archinal2018iau}.

\subsection{The internal-ranging observation model}
\label{sec:obsmodel}

We call this the \emph{internal-ranging} datum problem: the seven parameters of the
lunar frame are realized from range (distance) observations alone, with no
observation that directly fixes the frame's orientation, scale, or absolute
position. The Earth stations are externally known in the terrestrial frame, but a
range constrains only a distance, so the lunar-frame datum is bootstrapped from the
accumulated ranging geometry and its libration-driven diversity
(Theorem~\ref{thm:classification}). Throughout, \emph{radial} means along the
Earth--Moon line of sight, which for the near-side reflector cluster coincides with
the lunar body $+X$ (sub-Earth) axis; this coincidence is exactly why a lunocentric
scale change and an $X$-origin shift alias on near-side ranges
(Section~\ref{sec:radial}).

An LLR normal point measures the round-trip range between an Earth station at
inertial position $\bm{x}_{\mathrm{st}}(t)$ and a lunar retroreflector. A reflector
with PA body coordinates $\bm{p}$ appears in the inertial frame at
\begin{equation}
  \bm{r}(\bm{p},t) \;=\; \bm{x}_{\mathrm{M}}(t) + \Rmat(t)\,\bm{p},
  \label{eq:beacon}
\end{equation}
where $\bm{x}_{\mathrm{M}}(t)$ is the inertial lunar-centre position and
$\Rmat(t)$ the PA-body$\rightarrow$inertial orientation. We take both from the
JPL DE440 ephemeris \citep{park2021de440}, so that $\Rmat(t)$ carries the true
physical libration (Appendix~\ref{app:methods}); this distinction is essential and
is the mechanism behind item~(iii) of the classification theorem. The one-way
range and its line-of-sight unit vector are
\begin{equation}
  \varrho(\bm{p},t) = \bigl\|\bm{r}(\bm{p},t)-\bm{x}_{\mathrm{st}}(t)\bigr\|,
  \qquad
  \uhat(\bm{p},t) = \frac{\bm{r}(\bm{p},t)-\bm{x}_{\mathrm{st}}(t)}
                         {\varrho(\bm{p},t)}.
  \label{eq:range}
\end{equation}
A datum perturbation $\datum$ moves the reflector by $\Rmat(t)\bm{J}(\bm{p})\datum$;
projecting onto the sightline gives the per-observation datum sensitivity
\begin{equation}
  \bm{g}(\bm{p},t)^\top
  \;=\; \frac{\partial \varrho}{\partial \datum}
  \;=\; \uhat(\bm{p},t)^\top\,\Rmat(t)\,\bm{J}(\bm{p}) \;\in\;\R^{1\times7}.
  \label{eq:sensrow}
\end{equation}
Equation~\eqref{eq:sensrow} is the single object on which everything below rests:
it is a row vector, one per observation, and it is where the orientation $\Rmat(t)$
enters. The same construction gives the sensitivity of a two-station VLBI
differential delay and of an orbiter-to-beacon range once the appropriate
measurement gradient replaces $\uhat^\top$; we defer those to
Section~\ref{sec:multitechnique}.

\subsection{Fisher information and the degeneracy metric}
\label{sec:fisher}

For zero-mean Gaussian measurement noise with standard deviation $\sigma$, a set of
observations indexed by $i$ (each a station/beacon/epoch triple that passes a
visibility gate) yields the datum Fisher information matrix
\begin{equation}
  \Fisher \;=\; \sum_i \frac{1}{\sigma_i^2}\,
  \bm{g}_i\,\bm{g}_i^\top \;\in\;\R^{7\times7},
  \label{eq:fisher}
\end{equation}
the sum of rank-one outer products \citep{kay1993estimation}. The Cram\'er--Rao
bound $\mathrm{Cov}(\hat{\datum}) \succeq \Fisher^{-1}$ lower-bounds any unbiased
estimator; the datum \emph{defect} is $7-\operatorname{rank}(\Fisher)$, the number
of exactly unobservable directions.

We are interested not in the full seven-parameter covariance but in the
identifiability of the specific pair the ILRF struggles with:
$K=\{t_x,\,s\}=\{0,3\}$, marginalized over the remaining
$M=\{t_y,t_z,\theta_x,\theta_y,\theta_z\}=\{1,2,4,5,6\}$. The relevant object is
the Schur complement of the $K$ block after eliminating $M$,
\begin{equation}
  \Schur \;=\; \Fisher_{KK} - \Fisher_{KM}\,\Fisher_{MM}^{-1}\,\Fisher_{MK}
  \;\in\;\R^{2\times2},
  \label{eq:schur}
\end{equation}
the \emph{marginal} information on $(t_x,s)$ once all coupling to the other five
parameters has been accounted for \citep{zhang2005schur,boyd2004convex}. From
$\Schur$ we read three scalars, defined once and used throughout:
\begin{definition}[Degeneracy metric, origin CRLB, origin--scale correlation]
\label{def:metrics}
\begin{align}
  \mu &\;=\; \lmin(\Schur)
     && \text{(degeneracy metric)}, \label{eq:metric}\\
  c_x &\;=\; \sqrt{\bigl(\Schur^{-1}\bigr)_{00}}
     && \text{(origin-$X$ CRLB)}, \label{eq:crlb}\\
  \corr &\;=\; \frac{\bigl(\Schur^{-1}\bigr)_{01}}
                    {\sqrt{\bigl(\Schur^{-1}\bigr)_{00}\bigl(\Schur^{-1}\bigr)_{11}}}
     && \text{(origin--scale correlation)}. \label{eq:corr}
\end{align}
\end{definition}
The degeneracy metric $\mu\to0$ exactly when the pair becomes unobservable and
grows as the two directions separate; $\corr\to\pm1$ signals the near-degeneracy
that \citet{sosnica2025ilrf} observe. For the canonical case in which the
\emph{information} has the form
$\Schur=m\bigl[\begin{smallmatrix}1&\varrho\\\varrho&1\end{smallmatrix}\bigr]$
($m>0$), these reduce to closed forms ($\mu=\lmin(\Schur)=m(1-|\varrho|)$,
$c_x=m^{-1/2}(1-\varrho^2)^{-1/2}$, and $\corr=-\varrho$, so $|\varrho|=|\corr|$),
which we use as an analytic cross-check.

\paragraph{Preconditioning and units.}
The translation columns of $\bm{J}$ carry units of metres while the scale and
rotation columns carry units of the point coordinates ($\sim\!R_{\mathrm{Moon}}$),
giving $\Fisher$ a condition number $\sim\!10^{12}$. We precondition by dividing
columns $3$--$6$ by the mean lunar radius $R_{\mathrm{Moon}}=\SI{1737400}{m}$,
bringing all partials to $O(1)$ and the condition number to $\sim\!10^2$. This is a
positive diagonal congruence: it leaves the correlation $\corr$ (a pure
correlation coefficient) and the origin CRLB $c_x$ (which reads column $0$, left
unscaled) \emph{exactly invariant}, and it preserves rank and datum defect
(Appendix~\ref{app:precond}). The degeneracy metric $\mu$ is therefore reported in
these preconditioned, \emph{relative} units and is meaningful only for comparing
designs on a common footing---never as an absolute physical quantity. It is a
\Modelled{} figure of merit.
